% 117-06(117쪽) — 곡선 y=f(x)=x^2-5x+4 와 x축·y축·직선 x=k 로 생기는 세 부분 S_1, S_2, S_3. 스캔 화질이 낮아 TikZ 로 다시 그림.
% 원문에 있는 것만: 곡선 · 세 색칠한 부분과 라벨 S_1 S_2 S_3 · 눈금 4(y축) 1 4(x축) · 라벨 O, x, y, y=f(x), x=k.
% 직선 x=k 는 원문 그림처럼 4 보다 조금 오른쪽에 두었을 뿐 k 의 값은 원문에 없다(넣지 않는다).
\begin{tikzpicture}[x=0.50cm, y=0.30cm, line width=0.5pt]
  % S_1 (0 ≤ x ≤ 1)
  \fill[gray!18] plot[smooth, tension=0.6] coordinates {
    (0,4) (0.25,2.81) (0.50,1.75) (0.75,0.81) (1,0) } -- (0,0) -- cycle;
  % S_2 (1 ≤ x ≤ 4)
  \fill[gray!18] plot[smooth, tension=0.6] coordinates {
    (1,0) (1.25,-0.69) (1.50,-1.25) (1.75,-1.69) (2.00,-2) (2.25,-2.19) (2.50,-2.25)
    (2.75,-2.19) (3.00,-2) (3.25,-1.69) (3.50,-1.25) (3.75,-0.69) (4,0) } -- cycle;
  % S_3 (4 ≤ x ≤ k)
  \fill[gray!18] plot[smooth, tension=0.6] coordinates {
    (4,0) (4.25,0.81) (4.50,1.75) (4.75,2.81) (5.00,4) (5.25,5.31) (5.50,6.75) } -- (5.50,0) -- cycle;
  \draw[->, >=stealth, line width=0.7pt] (-0.95,0) -- (6.35,0) node[below=1pt, font=\small] {$x$};
  \draw[->, >=stealth, line width=0.7pt] (0,-3.10) -- (0,7.60) node[left=1pt, font=\small] {$y$};
  % 직선 x=k
  \draw[line width=0.55pt] (5.50,-1.20) -- (5.50,7.30);
  % 곡선
  \draw[line width=0.6pt] plot[smooth, tension=0.6] coordinates {
    (-0.45,6.45) (-0.30,5.59) (-0.15,4.77) (0,4) (0.25,2.81) (0.50,1.75) (0.75,0.81)
    (1,0) (1.25,-0.69) (1.50,-1.25) (1.75,-1.69) (2.00,-2) (2.25,-2.19) (2.50,-2.25)
    (2.75,-2.19) (3.00,-2) (3.25,-1.69) (3.50,-1.25) (3.75,-0.69) (4,0) (4.25,0.81)
    (4.50,1.75) (4.75,2.81) (5.00,4) (5.25,5.31) (5.55,7.05) };
  % S_1 안내선과 라벨
  \draw[line width=0.35pt] (0.52,0.95) .. controls (0.85,1.55) and (1.05,1.75) .. (1.30,1.90);
  \node[anchor=west, font=\small] at (1.30,2.05) {$S_1$};
  \node[font=\small] at (2.50,-0.95) {$S_2$};
  \node[font=\small] at (4.78,1.45) {$S_3$};
  % y=f(x) 안내선과 라벨
  \draw[line width=0.35pt] (4.42,6.25) .. controls (4.80,5.75) and (4.85,5.10) .. (5.05,4.35);
  \node[anchor=east, font=\small] at (4.40,6.55) {$y=f(x)$};
  % 눈금·라벨
  \node[left=2pt, font=\small] at (0,4) {$4$};
  \node[below left=0pt and -2pt, font=\small] at (0,0) {$\pt{O}$};
  \node[below=2pt, font=\small] at (1,0) {$1$};
  \node[below=2pt, font=\small] at (4,0) {$4$};
  \node[below=1pt, font=\small] at (5.50,-1.25) {$x=k$};
\end{tikzpicture}
